\chapter{Complexity — Classes}
%%% generated by the blueprint pipeline from TCSlib.Complexity.Randomized.Classes
\section{Overview}

Declarations of the Lean module \texttt{TCSlib.Complexity.Randomized.Classes}.

\section{Declarations}

\begin{definition}[Uniform probability of a predicate]
\lean{Randomized.randProb}
\label{Randomized.randProb}
\leanok
For a natural number $m$ and a decidable predicate $A$ on Boolean lists, let $X$ be
uniformly distributed over the $2^m$ Boolean lists of length $m$. The quantity defined
here is the probability that $A(X)$ holds.
\end{definition}

\begin{theorem}[Agreement preserves probability]
\lean{Randomized.randProb_congr}
\label{Randomized.randProb_congr}
\leanok
\uses{Randomized.randProb}
\difficulty{3}
Let $m\in\mathbb{N}$ and let $A$ and $B$ be predicates on Boolean words. If $A$ and $B$
agree on every Boolean word of length $m$, then their probabilities on a uniformly
random Boolean word of length $m$ are equal.
\end{theorem}

\begin{theorem}[Nonnegativity of random-string probability]
\lean{Randomized.randProb_nonneg}
\label{Randomized.randProb_nonneg}
\leanok
\uses{Randomized.randProb}
\difficulty{3}
For any nonnegative integer $m$ and any decidable property $A$ of Boolean strings, the
probability that a uniformly random Boolean string of length $m$ satisfies $A$ is
nonnegative.
\end{theorem}

\begin{theorem}[Number of Boolean Strings on a Finite Set]
\lean{Randomized.card_univ_bitstrings}
\label{Randomized.card_univ_bitstrings}
\leanok
\difficulty{3}
For every natural number $m$, the set of functions from an $m$-element set to
$\{\mathrm{false},\mathrm{true}\}$ has cardinality $2^m$.
\end{theorem}

\begin{theorem}[Probability is at most one]
\lean{Randomized.randProb_le_one}
\label{Randomized.randProb_le_one}
\leanok
\uses{Randomized.card_univ_bitstrings, Randomized.randProb}
\difficulty{4}
For any natural number $m$ and any predicate on finite Boolean words, the probability
that a uniformly random Boolean word of length $m$ satisfies the predicate is at most
$1$.
\end{theorem}

\begin{theorem}[Monotonicity of random-string probability]
\lean{Randomized.randProb_mono}
\label{Randomized.randProb_mono}
\leanok
\uses{Randomized.randProb}
\difficulty{3}
Let $m$ be a nonnegative integer, and let $A$ and $B$ be decidable properties of Boolean
strings. If every Boolean string of length $m$ satisfying $A$ also satisfies $B$, then
the probability that a uniformly random Boolean string of length $m$ satisfies $A$ is at
most the probability that it satisfies $B$.
\end{theorem}

\begin{theorem}[Probability of a complementary event]
\lean{Randomized.randProb_not}
\label{Randomized.randProb_not}
\leanok
\uses{Randomized.card_univ_bitstrings, Randomized.randProb}
\difficulty{5}
For any nonnegative integer $m$ and any decidable property $A$ of Boolean strings, the
probability that a uniformly random Boolean string of length $m$ does not satisfy $A$
equals one minus the probability that it satisfies $A$.
\end{theorem}

\begin{theorem}[Probability of a Certain Event]
\lean{Randomized.randProb_true}
\label{Randomized.randProb_true}
\leanok
\uses{Randomized.card_univ_bitstrings, Randomized.randProb}
\difficulty{3}
For every $m\in\mathbb{N}$, a uniformly random Boolean list of length $m$ satisfies the
property that is always true with probability $1$.
\end{theorem}

\begin{theorem}[Representing a List by a Function]
\lean{Randomized.exists_ofFn_eq}
\label{Randomized.exists_ofFn_eq}
\leanok
\difficulty{4}
Let $\ell$ be a bit string of length $m$. Then there is a function from the indices
$0,\ldots,m-1$ to $\{0,1\}$ whose associated list is exactly $\ell$.
\end{theorem}

\begin{theorem}[Product rule for independent word events]
\lean{Randomized.randProb_split}
\label{Randomized.randProb_split}
\leanok
\uses{Randomized.randProb}
\difficulty{6}
Let $m_1,m_2\in\mathbb N$, and let $A$ and $B$ be predicates on Boolean words. For a
uniformly random Boolean word of length $m_1+m_2$, the probability that its first $m_1$
bits satisfy $A$ and its remaining $m_2$ bits satisfy $B$ equals the product of the
probability that a uniformly random length-$m_1$ word satisfies $A$ and the probability
that a uniformly random length-$m_2$ word satisfies $B$.
\end{theorem}

\begin{theorem}[Probability of a Prefix Event]
\lean{Randomized.randProb_take}
\label{Randomized.randProb_take}
\leanok
\uses{Randomized.randProb, Randomized.randProb_congr, Randomized.randProb_split, Randomized.randProb_true}
\difficulty{4}
Let $m_1,m\in\mathbb{N}$ with $m_1\le m$, and let $A$ be any property of finite Boolean
lists. For a uniformly random Boolean list of length $m$, the probability that its first
$m_1$ entries satisfy $A$ equals the probability that a uniformly random Boolean list of
length $m_1$ satisfies $A$.
\end{theorem}

\begin{definition}[Polynomial length function]
\lean{Randomized.polyLen}
\label{Randomized.polyLen}
\leanok
For natural numbers $a$ and $k$, define a function $\mathbb{N} \to \mathbb{N}$ by
sending each natural number $n$ to $a(n+1)^k$.
\end{definition}

\begin{definition}[Verifier model]
\lean{Randomized.VerifierModel}
\statementsource{arora-barak-ch07-randomized-computation}{pdf-26d95080dbde-p003-b003}
\label{Randomized.VerifierModel}
\leanok
A verifier model consists of a predicate specifying which functions from pairs of
Boolean lists to optional Boolean outputs are efficient, and a predicate specifying
which functions from triples of Boolean lists to Boolean outputs are efficient.
\end{definition}

\begin{definition}[Boolean-valued verifier]
\lean{Randomized.boolVerifier}
\statementsource{arora-barak-ch07-randomized-computation}{pdf-26d95080dbde-p003-b003}
\label{Randomized.boolVerifier}
\leanok
Given a Boolean-valued verifier $M$ that takes an input list and a random-bit list, its
associated optional-output verifier returns the value of $M$ as a present Boolean output
on every input and random-bit list.
\end{definition}

\begin{definition}[Bounded-error probabilistic languages]
\lean{Randomized.InBPP}
\statementsource{Vad12}{pdf-bf2ce5422088-p023-b001}
\statementsource{arora-barak-ch07-randomized-computation}{
  pdf-26d95080dbde-p002-b003,
  pdf-26d95080dbde-p003-b003
}
\label{Randomized.InBPP}
\leanok
\uses{Randomized.boolVerifier, Randomized.polyLen, Randomized.randProb}
Fix a verifier model $E$ and a language $L$ of finite Boolean strings. The language $L$
is in bounded-error probabilistic polynomial time relative to $E$ if there are a
Boolean-valued verifier $M$ and natural numbers $a,k$ such that the verifier associated
to $M$ is efficient according to $E$, and for every input string $x$, membership in $L$
implies that $M$ outputs true with probability at least $2/3$, while nonmembership
implies that it outputs false with probability at least $2/3$. In both cases the
probability is taken over uniformly random Boolean strings of length $a(|x|+1)^k$,
as specified by the polynomial-length function in the definition.
\end{definition}

\begin{definition}[One-sided randomized language recognition]
\lean{Randomized.InRP}
\statementsource{Vad12}{
  pdf-bf2ce5422088-p021-b006,
  pdf-bf2ce5422088-p022-b001
}
\statementsource{arora-barak-ch07-randomized-computation}{
  pdf-26d95080dbde-p006-b002,
  pdf-26d95080dbde-p006-b003
}
\label{Randomized.InRP}
\leanok
\uses{Randomized.boolVerifier, Randomized.polyLen, Randomized.randProb}
Fix a verifier model $E$ and a language $L$ of finite Boolean lists. The language $L$
has an efficient verifier relative to $E$ if there exist a Boolean-valued verifier $M$
and natural numbers $a,k$ such that the optional-output verifier associated with $M$ is
efficient according to $E$, and for every Boolean list $x$, membership in $L$ implies
that $M$ accepts with probability at least $2/3$ when given a uniformly random string of
length $a(|x|+1)^k$, while nonmembership implies that $M$ rejects
for every random string of that length.
\end{definition}

\begin{definition}[Complementary one-sided probabilistic languages]
\lean{Randomized.InCoRP}
\statementsource{Vad12}{pdf-bf2ce5422088-p022-b003}
\statementsource{arora-barak-ch07-randomized-computation}{pdf-26d95080dbde-p006-b004}
\label{Randomized.InCoRP}
\leanok
\uses{Randomized.InRP}
Fix a verifier model $E$ and a language $L$ of Boolean lists. Then $L$ is in
complementary one-sided probabilistic polynomial time relative to $E$ exactly when the
complement of $L$ is in one-sided probabilistic polynomial time relative to $E$.
\end{definition}

\begin{definition}[Zero-error probabilistic languages]
\lean{Randomized.InZPP}
\statementsource{Vad12}{pdf-bf2ce5422088-p024-b001}
\statementsource{arora-barak-ch07-randomized-computation}{pdf-26d95080dbde-p007-b001}
\label{Randomized.InZPP}
\leanok
\uses{Randomized.polyLen, Randomized.randProb}
Fix a verifier model $E$ and a language $L$ of finite Boolean lists. The language $L$
belongs to zero-error probabilistic polynomial time relative to $E$ if there exist an
efficient verifier $M$, taking an input and a random Boolean list and returning either a
Boolean answer or no answer, and natural numbers $a,k$ such that, for every input $x$,
on every random list of length $a(|x|+1)^k$ the verifier never gives an incorrect
Boolean answer, and the probability that it gives no answer is at most $1/2$.
\end{definition}

\begin{definition}[Sequential race verifier]
\lean{Randomized.raceVerifier}
\label{Randomized.raceVerifier}
\leanok
Given two Boolean-output verifiers $M_1,M_2$ and a function $p_1$ from natural numbers
to natural numbers, define a verifier on input lists $x,r$ as follows: run $M_1$ on $x$
and the prefix of $r$ of length $p_1(|x|)$. If it returns true, output true; otherwise
run $M_2$ on $x$ and the suffix of $r$ after that prefix, outputting false if it returns
true and no answer otherwise.
\end{definition}

\begin{definition}[Strict majority verifier]
\lean{Randomized.majorityVerifier}
\label{Randomized.majorityVerifier}
\leanok
Given a Boolean-output verifier $M$ and functions $p,k$ from natural numbers to natural
numbers, define a verifier on Boolean lists $x,r$ by dividing $r$ into consecutive
blocks of length $p(|x|)$ and applying $M$ to $x$ and each of the first $k(|x|)$ blocks.
It outputs true exactly when strictly more than half of these $k(|x|)$ runs return true.
\end{definition}

\begin{definition}[Closure under racing verifiers]
\lean{Randomized.ClosedUnderRace}
\label{Randomized.ClosedUnderRace}
\leanok
\uses{Randomized.boolVerifier, Randomized.polyLen, Randomized.raceVerifier}
A verifier model $E$ is closed under racing if, for every pair of Boolean-output
verifiers $M_1,M_2$ and natural numbers $a,k$, whenever the Boolean verifiers associated
with $M_1$ and $M_2$ are efficient in $E$, the race verifier using $M_1,M_2$ and the
length function $n\mapsto a(n+1)^k$ is also efficient in $E$.
\end{definition}

\begin{definition}[Closure under answer predicates]
\lean{Randomized.ClosedUnderAnswerIs}
\label{Randomized.ClosedUnderAnswerIs}
\leanok
\uses{Randomized.boolVerifier}
A verifier model $E$ is closed under answer predicates if, for every verifier $M$ and
Boolean value $b$, whenever $M$ is efficient in $E$, the Boolean verifier that accepts
exactly when $M$ returns the answer $b$ is also efficient in $E$.
\end{definition}

\begin{definition}[Any-acceptance verifier]
\lean{Randomized.anyVerifier}
\label{Randomized.anyVerifier}
\leanok
Given a Boolean-output verifier $M$ and functions $p,k$ from natural numbers to natural
numbers, define a verifier on Boolean lists $x,r$ that outputs true exactly when at
least one of the first $k(|x|)$ consecutive blocks of $r$, each of length $p(|x|)$,
makes $M$ return true on input $x$.
\end{definition}

\begin{definition}[Closure under majority]
\lean{Randomized.ClosedUnderMajority}
\label{Randomized.ClosedUnderMajority}
\leanok
\uses{Randomized.boolVerifier, Randomized.majorityVerifier, Randomized.polyLen}
A verifier model $E$ is closed under majority if, for every Boolean-output verifier $M$
and natural numbers $a,k,a',k'$, whenever the Boolean verifier associated with $M$ is
efficient in $E$, the majority verifier using block length $n\mapsto a(n+1)^k$ and
repetition count $n\mapsto a'(n+1)^{k'}$ is also efficient in $E$.
\end{definition}

\begin{definition}[Closure under polynomial-length verifier combination]
\lean{Randomized.ClosedUnderAny}
\label{Randomized.ClosedUnderAny}
\leanok
\uses{Randomized.anyVerifier, Randomized.boolVerifier, Randomized.polyLen}
A verifier model $E$ is closed under the verifier combination if, for every verifier $M$
and all natural numbers $a,k,a',k'$, efficiency of the Boolean verifier associated with
$M$ implies efficiency of the Boolean verifier obtained from $M$ using the polynomial
lengths determined by $(a,k)$ and $(a',k')$.
\end{definition}

\begin{definition}[Efficiency closed under Boolean negation]
\lean{Randomized.ClosedUnderNot}
\statementsource{arora-barak-ch07-randomized-computation}{pdf-26d95080dbde-p013-b002}
\label{Randomized.ClosedUnderNot}
\leanok
\uses{Randomized.boolVerifier}
Fix a verifier model $E$. It is closed under Boolean negation if, for every
Boolean-valued verifier $M$, whenever the optional-output verifier associated with $M$
is efficient according to $E$, the verifier obtained by negating $M$'s output on every
input and random-bit list also has an efficient associated optional-output verifier
according to $E$.
\end{definition}

\begin{theorem}[Probability of an impossible event]
\lean{Randomized.randProb_false}
\label{Randomized.randProb_false}
\leanok
\uses{Randomized.randProb}
\difficulty{1}
For any verifier model, any natural number $m$, and a uniformly random Boolean word of
length $m$, the probability that the word satisfies an always-false condition is zero.
\end{theorem}

\begin{theorem}[Additivity for disjoint events]
\lean{Randomized.randProb_or_disjoint}
\label{Randomized.randProb_or_disjoint}
\leanok
\uses{Randomized.randProb}
\difficulty{4}
Let $m$ be a nonnegative integer, and let $A$ and $B$ be decidable properties of Boolean
strings. If no Boolean string of length $m$ satisfies both properties, then the
probability that a uniformly random Boolean string of length $m$ satisfies at least one
of them is the sum of their individual probabilities.
\end{theorem}

\begin{theorem}[Probability of a finite set of outputs]
\lean{Randomized.randProb_mem_eq_sum}
\label{Randomized.randProb_mem_eq_sum}
\leanok
\uses{Randomized.randProb, Randomized.randProb_congr, Randomized.randProb_false, Randomized.randProb_or_disjoint}
\difficulty{4}
Let $m$ be a nonnegative integer, let $g$ map Boolean strings to nonnegative integers,
and let $T$ be a finite set of nonnegative integers. Under the uniform distribution on
Boolean strings of length $m$, the probability that $g$ takes a value in $T$ is the sum,
over $j\in T$, of the probabilities that $g$ equals $j$.
\end{theorem}

\begin{theorem}[Complementary Boolean output probabilities]
\lean{Randomized.randProb_bool_false}
\label{Randomized.randProb_bool_false}
\leanok
\uses{Randomized.randProb, Randomized.randProb_congr, Randomized.randProb_not}
\difficulty{3}
Let $m\in\mathbb{N}$ and let $B$ map Boolean words to Boolean values. The probability
that $B$ returns false on a uniformly random Boolean word of length $m$ equals one minus
the probability that it returns true.
\end{theorem}

\begin{definition}[Count of accepted blocks]
\lean{Randomized.blockCount}
\label{Randomized.blockCount}
\leanok
For natural numbers $q$ and $K$, a Boolean-valued function $B$ on Boolean lists, and a
Boolean list $l$, define the block count as the number of integers $i$ with $0 \le i <
K$ such that $B$ evaluates to true on the list obtained by dropping the first $iq$
entries of $l$ and then taking the next $q$ entries.
\end{definition}

\begin{theorem}[Upper bound for a block count]
\lean{Randomized.blockCount_le}
\label{Randomized.blockCount_le}
\leanok
\uses{Randomized.blockCount}
\difficulty{3}
For natural numbers $q,K$, a predicate $B$ on finite bit strings, and a finite bit
string $l$, the number of the first $K$ consecutive length-$q$ blocks of $l$ satisfying
$B$ is at most $K$.
\end{theorem}

\begin{theorem}[Recursive formula for block counts]
\lean{Randomized.blockCount_succ}
\label{Randomized.blockCount_succ}
\leanok
\uses{Randomized.blockCount}
\difficulty{4}
For natural numbers $q,K$, a predicate $B$ on finite bit strings, and a finite bit
string $l$, the number of the first $K+1$ consecutive length-$q$ blocks of $l$
satisfying $B$ equals $1$ if the first length-$q$ prefix of $l$ satisfies $B$, and $0$
otherwise, plus the number of the first $K$ consecutive length-$q$ blocks of the string
obtained by dropping the first $q$ bits of $l$.
\end{theorem}

\begin{theorem}[Binomial Distribution of Successful Blocks]
\lean{Randomized.randProb_blockCount}
\label{Randomized.randProb_blockCount}
\leanok
\uses{Randomized.blockCount, Randomized.blockCount_le, Randomized.blockCount_succ, Randomized.randProb, Randomized.randProb_bool_false, Randomized.randProb_congr, Randomized.randProb_false, Randomized.randProb_or_disjoint, Randomized.randProb_split, Randomized.randProb_true}
\difficulty{6}
Fix a natural number $q$ and a Boolean-valued test $B$ on Boolean lists. For natural
numbers $K$ and $j$ with $j\le K$, divide a uniformly chosen Boolean string of length
$Kq$ into $K$ consecutive blocks of length $q$, and count the blocks on which $B$
returns true. If $s$ is the fraction of length-$q$ Boolean strings on which $B$ returns
true, then the probability that exactly $j$ blocks pass the test is $\binom Kj
s^j(1-s)^{K-j}$.
\end{theorem}

\begin{theorem}[Upper tail bound for successful blocks]
\lean{Randomized.randProb_tail_le}
\label{Randomized.randProb_tail_le}
\leanok
\uses{Randomized.blockCount, Randomized.blockCount_le, Randomized.randProb, Randomized.randProb_blockCount, Randomized.randProb_congr, Randomized.randProb_le_one, Randomized.randProb_mem_eq_sum, Randomized.randProb_nonneg}
\difficulty{6}
Let $q,K\in\mathbb N$, let $B$ assign a Boolean value to each Boolean word, and let
$\varepsilon\in\mathbb Q$ satisfy $0\le\varepsilon$. Suppose that for a uniformly random
Boolean word of length $q$, the probability that $B$ is true is at most
$1/2-\varepsilon$. For a uniformly random word of length $Kq$, divide it into $K$
consecutive blocks of length $q$, and let $N$ count the blocks on which $B$ is true.
Then the probability that $N\ge K/2$ is at most $2(1-4\varepsilon^2)^{\lfloor
K/2\rfloor}$.
\end{theorem}

\begin{theorem}[Exponential probability bound]
\lean{Randomized.one_sub_pow_le_half}
\label{Randomized.one_sub_pow_le_half}
\leanok
\difficulty{4}
Let $x$ be a rational number with $0 \le x \le 1$, and let $m$ be a natural number such
that $mx \ge 1$. Then $(1-x)^m \le 1/2$.
\end{theorem}

\begin{theorem}[Powered exponential bound]
\lean{Randomized.one_sub_pow_le_half_pow}
\label{Randomized.one_sub_pow_le_half_pow}
\leanok
\uses{Randomized.one_sub_pow_le_half}
\difficulty{3}
Let $x$ be a rational number with $0 \le x \le 1$, and let $m,T,e$ be natural numbers
such that $mx \ge 1$ and $mT \le e$. Then $(1-x)^e \le (1/2)^T$.
\end{theorem}

\begin{theorem}[Complement closure of bounded-error computation]
\lean{Randomized.InBPP.compl}
\label{Randomized.InBPP.compl}
\leanok
\uses{Randomized.ClosedUnderNot, Randomized.InBPP, Randomized.polyLen, Randomized.randProb, Randomized.randProb_congr}
\difficulty{4}
Let $E$ be a verifier model closed under Boolean negation: whenever it regards a Boolean
verifier as efficient, it also regards the verifier that negates every answer as
efficient. If a language $L$ belongs to the bounded-error class defined using $E$, then
its complement also belongs to that class.
\end{theorem}

\begin{theorem}[Converting a zero-error partial verifier to RP]
\lean{Randomized.inRP_of_zeroError}
\label{Randomized.inRP_of_zeroError}
\leanok
\uses{Randomized.ClosedUnderAnswerIs, Randomized.ClosedUnderAny, Randomized.InRP, Randomized.anyVerifier, Randomized.exists_ofFn_eq, Randomized.polyLen, Randomized.randProb, Randomized.randProb_congr, Randomized.randProb_mono, Randomized.randProb_nonneg, Randomized.randProb_not, Randomized.randProb_split, Randomized.randProb_take}
\difficulty{6}
Let $E$ be a verifier model with two closure properties: it can efficiently turn an
efficient partial verifier into the Boolean test that its output equals a fixed bit, and
it can efficiently take the specified disjunction of polynomially many tests on disjoint
random-tape blocks. Suppose $M$ is efficient, $a,k$ are natural numbers, $L'$ is a
language, and $b$ is a bit. For each $x\in L'$, assume that on every random tape of
length $a(|x|+1)^k$, $M$ never outputs the bit opposite to $b$ and the probability it
outputs no answer is at most $1/2$; for each $x\notin L'$, assume that on every such
tape it never outputs $b$. Then $L'$ has an efficient one-sided randomized verifier in
$E$, accepting members with probability at least $2/3$ and rejecting nonmembers with
certainty.
\end{theorem}

\begin{theorem}[Characterization of zero-error randomized computation]
\lean{Randomized.inZPP_iff_inRP_and_inCoRP}
\label{Randomized.inZPP_iff_inRP_and_inCoRP}
\leanok
\uses{Randomized.ClosedUnderAnswerIs, Randomized.ClosedUnderAny, Randomized.ClosedUnderRace, Randomized.InCoRP, Randomized.InRP, Randomized.InZPP, Randomized.anyVerifier, Randomized.boolVerifier, Randomized.exists_ofFn_eq, Randomized.inRP_of_zeroError, Randomized.polyLen, Randomized.raceVerifier, Randomized.randProb, Randomized.randProb_congr, Randomized.randProb_le_one, Randomized.randProb_nonneg, Randomized.randProb_not, Randomized.randProb_split, Randomized.randProb_take}
\difficulty{6}
Let $E$ be a verifier model closed under racing two efficient Boolean verifiers, testing
whether an efficient verifier returns a specified answer, and taking polynomially many
independent verifier trials. For every language $L$ of bit strings, $L$ has an efficient
polynomial-witness randomized verifier that never gives an incorrect answer and outputs
no answer with probability at most $1/2$ if and only if both $L$ and its complement have
efficient polynomial-witness randomized verifiers that accept members with probability
at least $2/3$ and never accept nonmembers.
\end{theorem}

\begin{definition}[Weak advantage parameter]
\lean{Randomized.weakAdv}
\label{Randomized.weakAdv}
\leanok
For natural numbers $c$ and $n$, define the weak advantage parameter to be the rational
number $\min\left\{\frac{1}{6},\,((n+1)^{-1})^c\right\}$.
\end{definition}

\begin{theorem}[Positivity of the weak advantage]
\lean{Randomized.weakAdv_pos}
\label{Randomized.weakAdv_pos}
\leanok
\uses{Randomized.weakAdv}
\difficulty{3}
For natural numbers $c$ and $n$, the quantity $\min\{1/6,((n+1)^{-1})^c\}$ is positive.
\end{theorem}

\begin{theorem}[Upper bound for weak advantage]
\lean{Randomized.weakAdv_le_sixth}
\label{Randomized.weakAdv_le_sixth}
\leanok
\uses{Randomized.weakAdv}
\difficulty{1}
For natural numbers $c$ and $n$, the weak advantage $\min\{1/6,((n+1)^{-1})^c\}$ is at
most $1/6$.
\end{theorem}

\begin{theorem}[Lower bound for weak advantage]
\lean{Randomized.weakAdv_ge}
\label{Randomized.weakAdv_ge}
\leanok
\uses{Randomized.weakAdv}
\difficulty{3}
For natural numbers $c$ and $n$, the quantity $\frac{((n+1)^{-1})^c}{6}$ is at most the
weak advantage $\min\{1/6,((n+1)^{-1})^c\}$.
\end{theorem}

\begin{theorem}[Weak advantage tail bound]
\lean{Randomized.weakAdv_tail_bound}
\label{Randomized.weakAdv_tail_bound}
\leanok
\uses{Randomized.one_sub_pow_le_half_pow, Randomized.polyLen, Randomized.weakAdv, Randomized.weakAdv_ge, Randomized.weakAdv_le_sixth, Randomized.weakAdv_pos}
\difficulty{5}
For natural numbers $c,d,n$, let $\varepsilon=\min\{1/6,((n+1)^{-1})^c\}$. Then
\[
2(1-4\varepsilon^2)^{\left\lfloor 56(n+1)^{2c+d}/2\right\rfloor}\le (1/2)^{(n+1)^d+1}.
\]
\end{theorem}

\begin{definition}[Weak randomized recognition of a language]
\lean{Randomized.InBPPWeak}
\statementsource{arora-barak-ch07-randomized-computation}{
  pdf-26d95080dbde-p007-b003,
  pdf-26d95080dbde-p007-b004
}
\label{Randomized.InBPPWeak}
\leanok
\uses{Randomized.boolVerifier, Randomized.polyLen, Randomized.randProb, Randomized.weakAdv}
Let $E$ be a verifier model, let $c$ be a natural number, and let $L$ be a language of
Boolean strings. The property holds if there are a Boolean-valued verifier $M$ on pairs
of Boolean strings and natural numbers $a,k$ such that the verifier induced by $M$ is
efficient with respect to $E$, and for every Boolean string $x$, both of the following
bounds hold: if $x\in L$, the probability that $M$ accepts using a random string of
length $\operatorname{polyLen}(a,k,|x|)$ is at least
$1/2+\operatorname{weakAdv}(c,|x|)$; if $x\notin L$, the probability that $M$ rejects
using a random string of that length is at least the same quantity.
\end{definition}

\begin{definition}[Strong bounded-error randomized recognition]
\lean{Randomized.InBPPStrong}
\statementsource{arora-barak-ch07-randomized-computation}{
  pdf-26d95080dbde-p002-b003,
  pdf-26d95080dbde-p003-b003
}
\label{Randomized.InBPPStrong}
\leanok
\uses{Randomized.boolVerifier, Randomized.polyLen, Randomized.randProb}
A language $L$ of Boolean strings has strong randomized polynomial-time recognition with
parameter $d$ if there are a Boolean-valued verifier $M$, natural numbers $a,k$, and an
efficient verifier induced by $M$, such that for every input $x$, if $x\in L$, then with
a polynomially bounded number of fair random bits the probability that $M$ accepts is at
least $1-(1/2)^{(|x|+1)^d+1}$; if $x\notin L$, the probability that it rejects is at
least the same bound.
\end{definition}

\begin{theorem}[Amplification of Weak Randomized Decision]
\lean{Randomized.bpp_error_reduction}
\label{Randomized.bpp_error_reduction}
\leanok
\uses{Randomized.ClosedUnderMajority, Randomized.InBPPStrong, Randomized.InBPPWeak, Randomized.blockCount, Randomized.majorityVerifier, Randomized.polyLen, Randomized.randProb, Randomized.randProb_bool_false, Randomized.randProb_congr, Randomized.randProb_mono, Randomized.randProb_not, Randomized.randProb_tail_le, Randomized.weakAdv, Randomized.weakAdv_pos, Randomized.weakAdv_tail_bound}
\difficulty{6}
Let $E$ be a verifier model closed under majority amplification: whenever a Boolean
verifier is efficient for $E$, so is the majority of polynomially many runs on
polynomial-length random blocks. If a language $L$ has an efficient verifier whose
acceptance probability on members and rejection probability on nonmembers are each at
least $1/2+\min\{1/6,(|x|+1)^{-c}\}$ for some natural number $c$, then for every natural
number $d$, $L$ has an efficient verifier whose corresponding probabilities are each at
least $1-2^{-((|x|+1)^d+1)}$.
\end{theorem}

\begin{theorem}[Strong randomized verification implies BPP]
\lean{Randomized.InBPPStrong.toInBPP}
\label{Randomized.InBPPStrong.toInBPP}
\leanok
\uses{Randomized.InBPP, Randomized.InBPPStrong}
\difficulty{3}
Let $E$ be a verifier model, let $d$ be a nonnegative integer, and let $L$ be a language
of bit strings. If there is an $E$-efficient randomized verifier using a polynomially
bounded number of random bits such that, on inputs of length $n$, it accepts members and
rejects nonmembers with probability at least $1-2^{-((n+1)^d+1)}$, then $L$ has an
$E$-efficient randomized verifier with both correctness probabilities at least $2/3$.
\end{theorem}

\begin{theorem}[Weak and standard BPP coincide]
\lean{Randomized.inBPPWeak_iff_inBPP}
\label{Randomized.inBPPWeak_iff_inBPP}
\leanok
\uses{Randomized.ClosedUnderMajority, Randomized.InBPP, Randomized.InBPPStrong.toInBPP, Randomized.InBPPWeak, Randomized.bpp_error_reduction, Randomized.weakAdv, Randomized.weakAdv_le_sixth}
\difficulty{3}
Let $E$ be a verifier model whose efficient verifiers remain efficient when amplified by
majority repetition, and let $c$ be a natural number. For every language $L$ of bit
strings, $L$ has a polynomial-time verifier with inverse-polynomial advantage over
random guessing, with the advantage specified by $c$, if and only if $L$ has a
polynomial-time verifier with acceptance probability at least $2/3$ on members and
rejection probability at least $2/3$ on nonmembers.
\end{theorem}
